\subsection{Closure-flow curvature regime}

\begin{proposition}[Directed joins are the Scott-continuity test carrier]\label{thm:directed-family-carrier}\lean{Infomaterial.Closure.Scott.DirectedFamily}\leanok\uses{thm:closure-regime-data, thm:kle}
A directed family is the finite charge-lattice projection of a directed subset.
It is the carrier on which a closure operator is tested for commuting with
directed joins.
\end{proposition}

\begin{theorem}[Finite charge closures are Scott-continuous]\label{thm:finite-scott-continuous}\lean{Infomaterial.Closure.Scott.finite_scott_continuous}\leanok\uses{thm:directed-family-carrier, thm:closure-operator, thm:spine-is-strict-regime}
On the finite charge carrier, every directed join is attained.  Therefore every
charge closure commutes with directed joins and sits on the strict/Scott side
of the regime projection.
\end{theorem}

\begin{proposition}[A non-Scott firing is failure to commute with a directed join]\label{thm:nsc-firing-carrier}\lean{Infomaterial.Flow.NSCFiring.IsNSC}\leanok\uses{thm:finite-scott-continuous}
An NSC firing is the precise obstruction equation: firing after the directed
join is not the directed join of the firings.  This is the incidence point
where infinitary antecedents enter the proof tree.
\end{proposition}

\begin{theorem}[Scott-continuity is absence of NSC firing]\label{thm:scott-continuous-iff-no-nsc-firing}\lean{Infomaterial.Flow.NSCFiring.scott_continuous_iff_no_nsc_firing}\leanok\uses{thm:nsc-firing-carrier, thm:finite-scott-continuous}
A charge closure is Scott-continuous exactly when no directed-family prefix is
an NSC firing for its operator.  This turns the regime flag into an incidence
condition.
\end{theorem}

\begin{theorem}[NSC firings source the non-Scott regimes]\label{thm:nsc-firing-source-of-aperiodic}\lean{Infomaterial.Flow.NSCFiring.nsc_firing_source_of_aperiodic}\leanok\uses{thm:scott-continuous-iff-no-nsc-firing, thm:regime-cohomology-vanishes-iff}
Genuine finite charge closures have no NSC firing, while a concrete
forall-antecedent firing operator witnesses one.  The witness lands on the
cohomology-bearing non-Scott side of the regime table.
\end{theorem}

\begin{theorem}[Curvature is the non-Scott regime projection]\label{thm:curvature-regime-correspondence}\lean{Infomaterial.Flow.CurvatureRegime.curvature_regime_correspondence}\leanok\uses{thm:nsc-firing-source-of-aperiodic, thm:regime-cohomology-vanishes-iff}
Flatness, vanishing \(H^1\), absence of curvature events, and strict/Scott
regime membership are the same side of the projection.  Curvedness,
nonvanishing \(H^1\), and NSC firing live on the transfinite-bounded or
aperiodic side.
\end{theorem}

\begin{theorem}[Positive curvature rank has an NSC firing witness]\label{thm:curvature-rank-has-nsc-firing}\lean{Infomaterial.Flow.CurvatureRegime.rank_nonzero_has_nsc_firing}\leanok\uses{thm:curvature-regime-correspondence}
The positive-rank side is not empty bookkeeping: the formalized
forall-antecedent operator supplies an actual NSC firing witness.
\end{theorem}
